\section{问题分析}
\subsection{问题一：单区组批}
本问的关键是将连续载荷能力转化为离散货箱批次。各机型对每个服务区的往返能力不同，且货箱不可拆分，故需先获得机型—服务区安全载荷，再在质量和容积限制下进行组合。结果要同时报告批次数和能耗，避免只按质量容量而忽略续航。

\subsection{问题二：时限排程}
组批确定后，问题转化为带时间窗的异构资源排程。每项运输任务占用一架无人机和一组对应电池；飞行结束后无人机返场，电池还需充电才能再次使用。双服务区路线可减少重复起降，但会延长任务时间并改变逐箱送达时刻，因此需将路线、起飞时刻和资源占用统一表示。

\subsection{问题三：通信保障}
通信配置继承运输轨迹。地面站—运输机直连可用时不额外安排中继；若直连无法支撑全程通信，则用运输机—中继机与中继机—地面站两跳链路覆盖缺口。中继任务包含转场与服务时间，并消耗中继机和能源组件。

\subsection{问题四：固定任务分区}
跨服务区运输架次形成不可拆分关系，可先构造服务区连接图并取连通分量，再将分量分配给独立任务组。分区之后，各组不能共享库存，资源总需求由组内并发占用的峰值决定；因此比较对象应为峰值并发需求相对库存的缺口，而不是简单累加各区任务数。
